\documentclass[12pt,openany]{book}

% ============= PACKAGES
\usepackage[utf8]{inputenc}
\usepackage[T1]{fontenc}
\usepackage[brazilian]{babel}

\usepackage{amsmath, amssymb, amsthm, mathtools}
\usepackage{geometry}
\usepackage{hyperref}
\usepackage{xcolor}
\usepackage[most]{tcolorbox}
\usepackage{titlesec}
\usepackage{fancyhdr}
\usepackage{graphicx}
\usepackage{tikz}
\usepackage{enumitem}
\usepackage{booktabs}
\usepackage{pgfplots}
\usepackage[portuguese,algoruled,lined,longend,linesnumbered]{algorithm2e}
\pgfplotsset{compat=1.18}

\geometry{margin=2.5cm}

% ============= COLORS
\definecolor{maincolor}{RGB}{0, 60, 120}
\definecolor{accentcolor}{RGB}{180, 30, 30}
\definecolor{lightcolor}{RGB}{230, 240, 255}
\definecolor{examplecolor}{RGB}{245, 245, 245}

% ============= LINKS
\hypersetup{
    colorlinks=true,
    linkcolor=maincolor,
    urlcolor=maincolor
}

% ============= PARAGRAPHS
\setlength{\parindent}{0pt}
\setlength{\parskip}{0.6em}

% ============= HELPERS
\newcommand{\R}{\mathbb{R}}
\newcommand{\N}{\mathbb{N}}
\newcommand{\Z}{\mathbb{Z}}
\newcommand{\Q}{\mathbb{Q}}
\newcommand{\eps}{\varepsilon}

\DeclareMathOperator{\im}{im}
\DeclareMathOperator{\dom}{dom}

\newcommand{\drawarray}[1]{
    \begin{tikzpicture}[
        cell/.style={
            draw,
            minimum width=1cm,
            minimum height=1cm
        }
    ]

    \foreach \v [count=\i from 1] in {#1} {

        \node[cell] (a\i) at (\i,0) {\v};

        \node[below=2pt] at (a\i.south) {\scriptsize \i};
    }

    \end{tikzpicture}
}

\SetKwFunction{Print}{Imprima}

\SetKwInput{KwInput}{Entrada}
\SetKwInput{KwOutput}{Saída}

\SetKwIF{If}{ElseIf}{Else}{se}{então}{senão se}{senão}{fim se}

\SetKwFor{For}{para}{faça}{fim}
\SetKw{KwTo}{até}
\SetKwFor{While}{enquanto}{faça}{fim}

\SetKw{Return}{retorne}

\SetKw{KwAnd}{e}
\SetKw{KwOr}{ou}
\SetKw{KwNot}{não}

\SetKw{Null}{null}
\SetKw{True}{verdadeiro}
\SetKw{False}{falso}

\SetKw{Continue}{continue}
\SetKw{Break}{break}

\newcommand{\complexity}[2]{
    \begin{complexitybox}
        \textbf{Tempo:} $O(#1)$ \par
        \textbf{Memória:} $O(#2)$
    \end{complexitybox}
}

% ============= COUNTERS
\newcounter{algorithmctr}

% ============= CHAPTER STYLE
\titleformat{\chapter}[display]
  {\Huge\bfseries\color{maincolor}}
  {\filleft\Huge\thechapter}
  {2ex}
  {\titlerule\vspace{1ex}\filright}
  [\vspace{1ex}\titlerule]

% ============= HEADER/FOOTER
\pagestyle{fancy}
\fancyhf{}

\fancyhead[L]{\nouppercase{\leftmark}}
\fancyhead[R]{$\sum$ymbolica}

\fancyfoot[C]{\thepage}

\renewcommand{\headrulewidth}{0.4pt}
\renewcommand{\footrulewidth}{0pt}

% ============= BOX (tcolorbox)
\tcbset{
    defstyle/.style={
        colback=lightcolor,
        colframe=maincolor,
        fonttitle=\bfseries,
        sharp corners
    },
    thmstyle/.style={
        colback=green!5,
        colframe=green!50!black,
        fonttitle=\bfseries,
        sharp corners
    },
    examplestyle/.style={
        colback=examplecolor,
        colframe=black!50,
        fonttitle=\bfseries,
        sharp corners
    },
    intuitionstyle/.style={
        colback=yellow!10,
        colframe=yellow!50!black,
        fonttitle=\bfseries,
        sharp corners
    },
    exercise/.style={
        colback=violet!5,
        colframe=violet!60!black,
        fonttitle=\bfseries,
        sharp corners
    }
}

% ============= ENVIRONMENTS
\newtcbtheorem[number within=chapter]{definition}{Definição}{defstyle}{def}
\newtcbtheorem[number within=chapter]{theorem}{Teorema}{thmstyle}{th}
\newtcbtheorem[number within=chapter]{lemma}{Lema}{thmstyle}{lem}
\newtcbtheorem[number within=chapter]{corollary}{Corolário}{thmstyle}{cor}
\newtcbtheorem[number within=chapter]{proposition}{Proposição}{thmstyle}{prop}
\newtcbtheorem[number within=chapter]{example}{Exemplo}{examplestyle}{ex}
\newtcbtheorem[number within=chapter]{remark}{Observação}{examplestyle}{rem}
\newtcbtheorem[number within=chapter]{exercise}{Exercício}{exercise}{exe}
\newtcolorbox{complexitybox}{
    colback=green!5,
    colframe=green!50!black,
    title=Complexidade,
    fonttitle=\bfseries
}
\newtcolorbox[
    auto counter,
    number within=section
]{algorithmbox}[2][]{
    enhanced,
    breakable,
    colback=blue!3,
    colframe=blue!60!black,
    fonttitle=\bfseries,
    title={Algoritmo~\thetcbcounter\;--\;#2},
    boxrule=0.8pt,
    arc=2mm,
    left=2mm,
    right=2mm,
    top=1mm,
    bottom=1mm,
    #1
}
\newtcolorbox{analysisbox}[1][]{
    enhanced,
    breakable,
    colback=yellow!5,
    colframe=yellow!50!black,
    fonttitle=\bfseries,
    title={Análise},
    boxrule=0.8pt,
    arc=2mm,
    left=2mm,
    right=2mm,
    top=1mm,
    bottom=1mm,
    #1
}
\newtcolorbox{intuitionbox}{
    intuitionstyle,
    title=Intuição
}
\newtcolorbox{taskbox}[1][]{
    enhanced,
    breakable,
    colback=red!3,
    colframe=red!55!black,
    fonttitle=\bfseries,
    title={Tarefa},
    boxrule=0.8pt,
    arc=2mm,
    left=2mm,
    right=2mm,
    top=1mm,
    bottom=1mm,
    #1
}

\renewcommand{\proofname}{Prova}

% ============= COVER
\newcommand{\makebookcover}{
\begin{titlepage}
    \begin{tikzpicture}[remember picture, overlay]
        \fill[maincolor] (current page.north west)
            rectangle ([xshift=4cm]current page.south west);
    \end{tikzpicture}

    \vspace*{3cm}

    \hspace{4.5cm}
    \begin{minipage}{0.6\textwidth}
        {\Huge\bfseries \textcolor{maincolor}{Algoritmos e Estruturas de Dados} \par}
        \vspace{0.5cm}
        {\Large Soluções Comentadas \par}

        \vspace{2cm}

        {\Large $\sum$ymbolica \par}

        \vfill

        {\large \today \par}
    \end{minipage}
\end{titlepage}
}

\usepackage{tikz}

% ============= DOCUMENT
\begin{document}

\makebookcover

\tableofcontents

\chapter{Lista 01}
\section{O maior Número Ímpar}

Dada uma lista $L$ contendo $n$ números inteiros.

Por exemplo:

\begin{center}
    \drawarray{9,42,21,14,28,3,19,32,46,6}
\end{center}

\begin{taskbox}
    Imprimir o maior número ímpar armazenado na lista ou, caso não haja, imprimir "Não há ímpar!".
\end{taskbox}

No exemplo acima, isso nos daria
\[
    21
\]

\subsection*{Soluções}

A ideia principal do algoritmo é fazer uma varredura linear na lista armazenando sempre o maior número ímpar encontrado até o momento numa variável $m$.
Note que a inicialização de $m$ é fundamental, uma vez que não sabemos se há (e onde estaria) o primeiro número ímpar.
Nesse sentido, podemos tomar um valor para $m$ que indique que não foi encontrado nenhum número ímpar, como, por exemplo, $m = 0$.
Após percorrer toda a lista atualizando o valor de $m$ quando encontrado o primeiro ímpar ou quando encontrado algum valor $L[i] > m$, no final teremos ou $m = 0$ e portanto não há nenhum ímpar na lista, ou $m = L[i]$ para algum $1 \leq i \leq n$.

\begin{algorithm}[H]
    \DontPrintSemicolon
    
    \KwInput{Uma lista $L$ com $n$ números inteiros}

    $m \gets 0$
    
    \For{$i \gets 1$ \KwTo $n$}{
        \If{$L[i]$ é ímpar \KwAnd ($m = 0$ \KwOr $L[i] > m$)}{
            $m \gets L[i]$
        }
    }

    \BlankLine

    \If{$m = 0$}{
        \tcc{Não há nenhum número ímpar na lista}
        \Print{"Não há ímpar!"}
    }
    \Else{
        \Print{$m$}
    }
\end{algorithm}

\complexity{n}{1}

\section{O Segundo Maior Número Ímpar}
Dada uma lista $L$ contendo $n$ números inteiros.

Por exemplo:

\begin{center}
    \drawarray{9,42,21,14,28,3,19,32,46,6}
\end{center}

\begin{taskbox}
    Imprimir o segundo maior número ímpar armazenado na lista ou, caso não haja, imprimir "Não existe!".
\end{taskbox}

No exemplo acima, isso nos daria
\[
    19
\]

\subsection*{Soluções}
\subsubsection*{Varredura Dupla}
A ideia principal do algoritmo é fazer uma varredura linear na lista armazenando sempre o maior número ímpar encontrado até o momento numa variável $m_1$.
Note que a inicialização de $m_1$ é fundamental, uma vez que não sabemos se há (e onde estaria) o primeiro número ímpar.
Nesse sentido, podemos tomar um valor para $m$ que indique que não foi encontrado nenhum número ímpar, como, por exemplo, $m_1 = 0$.
Após percorrer toda a lista atualizando o valor de $m_1$ quando encontrado o primeiro ímpar ou quando encontrado algum valor $L[i] > m_1$, no final teremos ou $m_1 = 0$ e portanto não há nenhum ímpar na lista, ou $m_1 = L[i]$ para algum $1 \leq i \leq n$.
Uma vez armazenado o maior número ímpar (se existir, caso não exista, basta retornar), podemos fazer uma outra varredura linear para o segundo maior número ímpar $m_2$, agora verificando, também, se $L[i] < m_1$.

\begin{algorithm}[H]
    \DontPrintSemicolon
    
    \KwInput{Uma lista $L$ com $n$ números inteiros}

    \tcc{Busca pelo maior ímpar}

    $m_1 \gets 0$
    
    \For{$i \gets 1$ \KwTo $n$}{
        \If{$L[i]$ é ímpar \KwAnd ($m_1 = 0$ \KwOr $L[i] > m_1$)}{
            $m_1 \gets L[i]$
        }
    }

    \BlankLine

    \If{$m_1 = 0$}{
        \tcc{Não há nenhum número ímpar na lista}
        \Print{"Não há ímpar!"} \\
        \Return
    }

    \BlankLine

    \tcc{Busca pelo segundo maior ímpar}
    
    $m_2 \gets 0$

    \BlankLine

    \For{$i \gets 1$ \KwTo $n$}{
        \If{$L[i]$ é ímpar \KwAnd $L[i] < m_1$ \KwAnd ($m_2 = 0$ \KwOr $L[i] > m_2$)}{
            $m_2 \gets L[i]$
        }
    }
    
    \BlankLine

    \If{$m_2 = 0$}{
        \Print{"Não existe!"}
    }
    \Else {
        \Print{$m_2$}
    }
\end{algorithm}

\subsubsection*{Varredura Única - Dupla Verificação}
Uma outra possibilidade é fazer duas verificações ao invés de uma.
Nesse caso, sempre que encontrarmos um número ímpar checamos primeiramente se ele deve ir para $m_1$ e, caso afirmativo,
movemos o conteúdo de $m_1$ para $m_2$, caso contrário, checamos se ele deve ir para $m_2$.

\begin{algorithm}[H]
    \DontPrintSemicolon
    
    \KwInput{Uma lista $L$ com $n$ números inteiros}

    \tcc{Busca pelo maior ímpar}

    $m_1 \gets 0$ \\
    $m_2 \gets 0$
    
    \For{$i \gets 1$ \KwTo $n$}{
        \If{$L[i]$ é par}{
            \Continue
        }

        \BlankLine

        \If{$m_1 = 0$ \KwOr $L[i] > m_1$}{
            $m_2 \gets m_1$ \\
            $m_1 \gets L[i]$
        }
        \ElseIf{$m_2 = 0$ \KwOr $L[i] > m_2$}{
            $m_2 \gets L[i]$
        }
    }

    \BlankLine

    \If{$m_2 = 0$}{
        \Print{"Não existe!"}
    }
    \Else {
        \Print{$m_2$}
    }
\end{algorithm}

\complexity{n}{1}


\section{Busca Aproximada}
Dada uma lista $L$ contendo $n$ números inteiros.

Por exemplo:

\begin{center}
    \drawarray{9,42,21,14,25,3,19,33,46,6}
\end{center}

e um número inteiro $k$

\begin{taskbox}
    Encontrar o número $k$ na lista e imprimir "Encontrei!" ou, caso não esteja presente, encontrar o número mais próximo de $k$ e imprimi-lo.
\end{taskbox}

No exemplo acima, para $k = 31$, isso nos daria
\[
    33
\]

\subsection*{Soluções}
\subsubsection*{Varredura Linear}
Podemos varrer a lista mantendo salvo um índice $j$ (o índice com o valor mais próximo de $k$) e um valor $d_m$ (a diferença absoluta entre o valor mais próximo de $k$ e $k$) de forma que analisemos $d = \text{abs}(L[i] - k)$.
Caso $d = 0$, então $L[i] = k$ e podemos imprimí-lo e retornar.
Caso contrário, verifiquemos se $d_m < \text{abs}(L[j] - L[i])$, em caso positivo, salvamos $j = i$ e $d_m = \text{abs}(L[j] - L[i])$.

\begin{algorithm}[H]
    \DontPrintSemicolon
    
    \KwInput{Uma lista $L$ com $n$ números inteiros e um número inteiro $k$}

    $j \gets 0$ \\
    $d_m \gets 0$ \\

    \For{$i \gets 1$ \KwTo $n$}{
        $d \gets L[i] - k$ \\
        \If{$d = 0$}{
            \Print{"Encontrei!"} \\
            \Return
        }
        \ElseIf{$j = 0$ \KwOr $d_m < \text{abs}(L[j] - L[i])$}{
            $j \gets i$ \\
            $d_m \gets \text{abs}(L[j] - L[i])$
        }
    }

    \Print{$L[j]$}
\end{algorithm}

\complexity{n}{1}


\section{Ímpar-Ímpar}
Dada uma lista $L$ contendo $n$ números inteiros.

Por exemplo:

\begin{center}
    \drawarray{9,2,7,7,2,2,1,7,7,9}
\end{center}

\begin{taskbox}
    Verificar se existe algum número ímpar que aparece um número ímpar de vezes na lista imprimindo "Sim", caso haja e "Não" caso contrário.
\end{taskbox}

No exemplo acima, isso nos daria
\[
    1
\]

\subsection*{Soluções}
\subsubsection*{Análise Quadrática}
Podemos percorrer a lista usando dois ponteiros, $i$ e $j$.
Para cada $1 \leq i \leq n$ tal que $L[i]$ seja ímpar percorremos a lista no intervalo $i + 1 \leq j \leq n$ contando quantos elementos são iguais a $L[i]$ numa variável $c$.
Caso $c$ seja ímpar, encontramos um número ímpar que satisfaz o problema e podemos imprimir e retornar.
Caso não encontrado nenhum número, imprimimos que não há.

\begin{algorithm}[H]
    \DontPrintSemicolon
    
    \KwInput{Uma lista $L$ com $n$ números inteiros}

    \For{$i \gets 1$ \KwTo $n$}{
        $c \gets 1$ \\
        \For{$j \gets i + 1$ \KwTo $n$}{
            \If{$L[j] = L[i]$}{
                $c \gets c + 1$
            }
        }

        \BlankLine

        \If{$c$ é ímpar}{
            \Print{"Sim"} \\
            \Return
        }
    }

    \Print{"Não"}
\end{algorithm}

\begin{analysisbox}
    Note que, para cada $1 \leq i \leq n$ percorremos o intervalo $i + 1 \leq j \leq n$.
    Isso nos dá
    \[
        \begin{aligned}
            \sum_{i = 1}^{n} (n - (i + 1)) &= \sum_{i = 1}^{n} (n - i) - \sum_{i = 1}^{n} 1 \\
            &= \sum_{i = 1}^{n} i - n \\
            &= \frac{n(n + 1)}{2} - n \\
            & = \frac{n^2 - n}{2}
        \end{aligned}
    \]
    Logo,
    \[
        O(n^2)
    \]
\end{analysisbox}

\complexity{n^2}{1}

\subsubsection*{Algoritmo Misterioso}
\begin{exercise}{Custo $O(n)$}{}
    Deixamos como exercício para o leitor a construção de um algoritmo cujo custo assintótico amortizado de tempo é $O(n)$.

    \vspace{1em}
    
    \textbf{Observação:} não há restrição quanto ao uso de memória.
    Mas você pode impor limitações, por exemplo, $\forall 1 \leq i \leq n$, $1 \leq L[i] \leq n$, em suas construções para construções simplificadas.
\end{exercise}

\end{document}